% HAND-WRITTEN fragment -- NOT generated by maestro2tex. Input directly from
% AnalogChapter.tex, after the block sections: it refers to all of them by \ref.
%
% Replaces the six-fragment calibration section of 2026-08-21
% (note_calibration_{intro,paths,seq,residual}.tex + tab_calibration_errors.tex),
% which the owner judged far too long. The per-block measurement recipes moved
% into the block sections themselves as note_meas_*.tex; what remains here is
% only what firmware programs. Deleted fragments are recoverable from git.
%
% Sources of the numbers below, so they can be re-checked:
%   bias trim span / spread   tab_biasgen_trim; MC sigma/mu 12.5 % (bg_mc, Plan.4.Run.2)
%   A1 offset                 tab_ampab_mc_offset, tab_ampab_mc_tracking
%   zero-current offset       tab_px_mc_izero, sss:pixel-mc
%   reference droop           note_dacr2r12_supply (R^2 = 94 %, residual 931 uV rms
%                             against 25.70 mV pp), tab_dacr2r12_mc_reference
%   ladder INL                tab_dacr2r12_mc_ideal (1.49 DAC LSB mean = 0.788 mV)
%   converter step            tab_adc_transfer (1.591 mV effective LSB)
%   Rf tolerance / matching   tab_tiarprog_corners, tab_tiarprog_mc
%   quantisation floor        1.591/sqrt(12) = 0.459 mV; two conversions = 0.650 mV
%   potential residual        RSS(0.650, 0.931, 0.788) = 1.38 mV rms
%   current residual          RSS(0.650, 0.294) = 0.713 mV rms = 0.45 LSB on the
%                             measured 1.591 mV step. Published as 0.71 LSB until
%                             2026-08-25: the RSS is of two VOLTAGES (mV), and the
%                             unit was asserted, not converted.
%   range condition           0.815 V / (3 x 5.7176 mV) = 47.5
%   uncalibrated potential sigma  5.7176 mV (tab_ampab_mc_tracking)
% Library cell names are deliberately kept OUT of the rendered prose (owner
% request 2026-08-21); they live in these comments instead.

\subsection{Per-Chip Calibration}\label{ss:calibration}

Every correction is a stored digital constant except one, the bias generator's
6-bit trim code. Table~\ref{tab:calibration-constants} is the complete set of
knobs: what each corrects, how it is measured, and when it is applied. The
measurement itself is described in each block's section, under \emph{Measuring
this block on chip}.

Only two constants are absolute --- one voltage at the reference DAC output, one
current forced at the working electrode --- and both need the production tester.
The other eight are zeros and ratios the channel measures against itself through
its own converter, which is what lets the power-up refresh run with no equipment.

Order matters twice. At test, apply the bias trim before anything else is
measured: it moves every transconductance behind every constant that follows. In
a reading, subtract the zero code, scale by the gain constant for the active gain
code, then subtract the drop across the working-electrode switch and, where an
impedance sweep has supplied it, the uncompensated electrolyte resistance.
Potential-axis corrections apply where the waveform is tabulated, not per sample.

\input{\MaestroRoot tab_calibration_constants.tex}

\paragraph{Residual.} Each zero is the difference of two conversions, so its
floor is \SI{0.65}{\milli\volt} rms, two quantisations of the converter's
\SI{1.591}{\milli\volt} step; the common-mode DAC has to dither the input to
reach it. On the potential axis the residual is \SI{1.38}{\milli\volt} rms ---
that floor, the \SI{0.93}{\milli\volt} droop-fit residual and
\SI{0.79}{\milli\volt} of uncorrected ladder INL in quadrature --- against
\SI{5.71}{\milli\volt} of uncalibrated \(\sigma\). That holds over the
specified \SIrange{0.4}{2.1}{\volt} window only: outside it the tracking error
is no longer a constant and no stored constant removes it
(Section~\ref{sss:pixtop-mc}). On the current axis it is
\SI{0.71}{\milli\volt} rms --- \SI{0.45}{LSB} on the converter's measured
\SI{1.591}{\milli\volt} step (Table~\ref{tab:adc-transfer}) --- plus
\SIrange{0.1}{0.4}{\percent} of reading, against \SIrange{12.1}{167.3}{LSB}
uncalibrated, which are on the nominal \SI{2.441}{\milli\volt} step.
A constant removes the current zero only
while that zero stays inside the converter window, which requires
\(R_f/R_{\mathrm{cell}} < 47.5\); the \SI{560}{\kilo\ohm} tap into a
\SI{10}{\kilo\ohm} electrode (Section~\ref{sss:pixel-mc}) is outside it and those
samples are recoverable by nothing.

\paragraph{Out of reach.} The reference DAC's \(-6.6\)~LSB of DNL is lost command
resolution, up to \SI{3.5}{\milli\volt} at the worst code, and no constant
creates a code the ladder does not produce. A reference-electrode junction
potential is indistinguishable from A1's offset, so only a zero taken \emph{wet}
against that electrode absorbs it. The temperature dependence of every constant
here is unquantified: no block in this chapter has a temperature-resolved Monte
Carlo, so how far a power-up re-measurement tracks a \SI{37}{\celsius} operating
point is not known.
